\documentclass{rvcepaper}
\begin{document}

% ---------- Entrepreneurship & IPR (CIE - I)
\begin{iempaperB}
\paperinfo{ayear={2024-2025 (ODD Sem)},date={25\textsuperscript{th} November 2024},code={IM254TA},sem={V Semester},exam={CIE -- I},maxmarks={10 + 50},duration={120 Mins},
 title={ENTREPRENEURSHIP AND INTELLECTUAL PROPERTY RIGHTS}}
\iemheadB
\ptitle{PART -- A}
\begin{cietable}{M}{BT}{CO}
\q{1}{Define Compulsory license under the patent law.}{02}{1}{4}
\q{2}{List any two characteristics of successful entrepreneurs.}{02}{2}{1}
\q{3}{Difference between Intrapreneur \& Entrepreneur.}{02}{1}{1}
\q{4}{List any two non-patentable inventions under the Indian patent Act.}{02}{2}{2}
\q{5}{List any two myths of entrepreneurs.}{02}{2}{1}
\end{cietable}
\ptitle{PART -- B}
\begin{cietable}{M}{BT}{CO}
\q{1.a}{A software company develops an AI tool for legal research and applies for a patent in India. Discuss the challenges it might face under Indian patent law.}{05}{3}{3}
\q{b}{Explain the difference between a product patent and a process patent with an examples.}{05}{1}{4}
\q{2}{Create a commercialization strategy for a patent related to renewable energy technology.}{10}{3}{3}
\q{3}{Propose a framework for enhancing Traditional Knowledge protection under international Intellectual Property laws.}{10}{2}{4}
\q{4}{Create a case study on how perseverance as a characteristic helped a famous entrepreneur achieve success.}{10}{3}{1}
\q{5}{Explain how entrepreneurship fosters innovation in an economy with one case example}{10}{2}{1}
\end{cietable}
\btfoot{|c|c|c|*{11}{C{0.82cm}|}}{\hline
\multirow{3}{*}{\bfseries Marks Distribution} & \multirow{2}{*}{\bfseries\shortstack{Max\\Marks}} & \multicolumn{1}{c|}{\bfseries Particulars} & CO1 & CO2 & CO3 & CO4 & CO5 & L1 & L2 & L3 & L4 & L5 & L6\\\cline{3-14}
 & & Quiz & 04 & 02 & -- & 04 & -- & 02 & 02 & 06 & -- & -- & --\\\cline{3-14}
 & & Test & 20 & 10 & 15 & 15 & -- & 05 & 10 & 25 & -- & -- & --\\\hline}
\stars{*****}
\end{iempaperB}

% ---------- IM352IA Operations Management (CIE - I)
\begin{iempaperB}
\paperinfo{ayear={2024-2025 (ODD Sem)},date={25\textsuperscript{th} November 2024},code={IM352IA},sem={V},exam={CIE -- I},maxmarks={10 + 50},duration={120 Min},
 title={OPERATIONS MANAGEMENT}}
\iemheadB
\begin{cietable}{M}{BT}{CO}
\qpart{Part -- A}
\q{1}{Operations ensure the efficient use of \blank[1.8cm] to meet customer demands.}{01}{1}{1}
\q{2}{A supply chain includes the flow of materials, information, and \blank[1.8cm] across organizations.}{01}{1}{1}
\q{3}{Effective supply chain management ensures timely delivery and reduces \blank[1.8cm].}{01}{1}{1}
\q{4}{Competitive priorities like cost, quality, and \blank[1.8cm] determine how a company competes in the market.}{01}{1}{1}
\q{5}{The Fourth Industrial Revolution integrates technologies like IoT, \blank[1.8cm], and robotics into operations.}{01}{1}{1}
\q{6}{Processes can be categorized as \blank[1.8cm] or supporting based on their direct impact on the final product.}{01}{1}{1}
\q{7}{Differentiate between order winner and order qualifier.}{02}{2}{2}
\q{8}{A hospital is considering a new procedure to be offered at Rs.200 per patient. The fixed cost per year would be Rs.100,000, with total variable costs of Rs.100 per patient. What is the break-even quantity for this service?}{02}{2}{2}
\qpart{Part --B}
\q{1}{Define operations management. Briefly explain the key terms involved in the definition of operations management with a schematic model of a production system.}{10}{2}{1}
\q{2}{With suitable example, Distinguish between Manufacturing and service processes.}{10}{3}{1}
\q{3}{White Valley Ski Resort is planning the ski lift operation for its new ski resort. Management is trying to determine whether one or two lifts will be necessary; each lift can accommodate 250 people per day. Skiing normally occurs in the 14-week period from December to April, during which the lift will operate 7 days per week. The first lift will operate at 90 percent capacity if economic conditions are bad, the probability of which is believed to be about a 0.3. During normal times the first lift will be utilized at 100 percent capacity, and the excess crowd will provide 50 percent utilization of the second lift. The probability of normal times is 0.5. Finally, if times are really good, the probability of which is 0.2, the utilization of the second lift will increase to 90 percent. The equivalent annual cost of installing a new lift, recognizing the time value of money and the lift's economic life, is \$50,000. The annual cost of installing two lifts is only \$90,000 if both are purchased at the same time. If used at all, each lift costs \$200,000 to operate, no matter how low or high its utilization rate. Lift tickets cost \$20 per customer per day. Should the resort purchase one lift or two?}{20}{4}{3}
\q{4}{A manager is deciding whether to build a small or a large facility. Much depends on the future demand that the facility must serve, and demand may be small or large. The manager knows with certainty the payoffs that will result under each alternative, shown in the following payoff table.\newline
What is the best alternative for each decision rule: Maximin. Maximax. Laplace. Minimax Regret.
\qtabc{|l|c|c|}{\hline & \multicolumn{2}{c|}{Possible future demand}\\\hline Alternative & Low & High\\\hline Small facility & 200 & 270\\\hline Large facility & 160 & 800\\\hline Do nothing & 0 & 0\\\hline}}{10}{4}{3}
\end{cietable}
\end{iempaperB}

% ---------- IM3531A Quality Assurance (CIE - I)
\begin{iempaperB}
\paperinfo{ayear={2024-2025 (ODD Sem)},date={26\textsuperscript{th} November 2024},code={IM3531A},sem={V Semester},exam={CIE -- I},maxmarks={10 + 50},duration={120 Mins},
 title={QUALITY ASSURANCE}}
\iemheadB
\ptitle{PART -- A}
\begin{cietable}{M}{BT}{CO}
\q{1}{Differentiate between defect and defective.}{2}{2}{1}
\q{2}{What is the primary goal of quality assurance?}{1}{2}{1}
\q{3}{How does ISO certification benefit organizations?}{1}{1}{2}
\q{4}{Mention four key core technologies that drive Quality 4.0.}{2}{2}{3}
\q{5}{List the 7 QC tools}{2}{1}{3}
\q{6}{What is ATS and ARL?}{2}{1}{2}
\end{cietable}
\ptitle{PART -- B}
\begin{cietable}{M}{BT}{CO}
\q{1.}{A certain type of storage of battery lasts on an Average 3 years and Standard of 0.5, assume normal distribution, find percentage of battery that lasts a) Less than 2 years, b) Less than 4 years, c)between 2.5 to 4.5 years, d)less than 4 years.}{10}{3}{3}
\q{2.}{Describe the principles of Rational Subgrouping and apply them to a specific industry, such as manufacturing or service, to show how subgrouping can be used to reduce variation and improve process quality.}{10}{2}{1}
\q{3.}{Classify the Cost of Good Quality and Cost of Bad Quality in the automotive industry. Provide examples of each cost category and explain how they can be managed effectively to improve product performance and reduce overall expenses.}{10}{1}{2}
\q{4.}{Choose a specific product (e.g., a smartphone, automobile, or any consumer electronics) and analyze how the Eight Dimensions of Quality influence its acceptance by customers. Break down each dimension---such as performance, features, reliability, and serviceability---and evaluate their relative importance in shaping customer perceptions. How do these dimensions interact with each other to create a competitive advantage for the product in the market? Use relevant examples to support your analysis.}{10}{4}{4}
\q{5.}{What are the key benefits of implementing the ISO 14000 standard for an organization?}{10}{2}{2}
\end{cietable}
\btfoot{|c|c|c|*{11}{C{0.82cm}|}}{\hline
\multirow{3}{*}{\bfseries Marks Distribution} & \multicolumn{2}{c|}{\bfseries Particulars} & CO1 & CO2 & CO3 & CO4 & CO5 & L1 & L2 & L3 & L4 & L5 & L6\\\cline{2-14}
 & \multirow{2}{*}{\shortstack{Max\\Marks}} & Quiz & 03 & 06 & 04 & -- & -- & 05 & 05 & -- & -- & -- & --\\\cline{3-14}
 & & Test & 10 & 20 & 10 & 10 & -- & 10 & 20 & 10 & 10 & -- & --\\\hline}
\stars{*****}
\end{iempaperB}

% ---------- IM254TA Financial Accounting and Costing (CIE - I)
\begin{iempaperB}
\paperinfo{ayear={2024-2025 (ODD Sem)},date={26\textsuperscript{th} November 2024},code={IM254TA},sem={V Semester},exam={CIE -- I},maxmarks={10 + 50},duration={120 Mins},
 title={FINANCIAL ACCOUNTING AND COSTING}}
\iemheadB
\ptitle{PART -- A}
\begin{cietable}{M}{BT}{CO}
\q{1}{Write a note on double-entry bookkeeping.}{02}{2}{1}
\q{2}{List the main financial statements prepared in financial accounting}{02}{1}{1}
\q{3}{Mention the accounting equation}{02}{1}{1}
\q{4}{What is the purpose of a trial balance?}{02}{2}{1}
\q{5}{Depict the journal format in accounting.}{02}{1}{1}
\end{cietable}
\ptitle{PART -- B}
\begin{cietable}{M}{BT}{CO}
\q{1.}{Explain the steps involved in the accounting cycle and discuss their significance in ensuring accurate financial reporting}{10}{3}{1}
\q{2.}{Discuss how zero-based budgeting differs from traditional budgeting approaches. What are the benefits and limitations of adopting a zero-based budgeting system?}{10}{2}{1}
\q{3 a.}{Distinguish between Single Entry vs. Double Entry accounting process.}{05}{2}{1}
\q{b.}{Write a note on International Financial Reporting Standards.}{05}{2}{1}
\q{4 a.}{Journalise the business transactions during the month of Jan 2024.\newline
Jan 1: Owner invested \rupee 50,000 cash in the business.\newline
Jan 3: Purchased office furniture for \rupee 10,000, paid in cash.\newline
Jan 5: Bought goods worth \rupee 20,000 on credit from M/s XYZ Traders.\newline
Jan 10: Paid \rupee 5,000 to M/s XYZ Traders as partial payment.\newline
Jan 15: Sold goods for \rupee 15,000 on cash.\newline
Jan 20: Paid office rent of \rupee 3,000 in cash.}{05}{3}{2}
\q{b.}{Discuss the principles of debit and credit in journalizing transactions. How do these principles apply to different types of accounts?}{05}{3}{1}
\q{5.}{Journalise the business transactions during the month of Feb XXXX\newline
Feb 2: Started business with cash \rupee 1,00,000 and machinery \rupee 50,000.\newline
Feb 5: Opened a bank account by depositing \rupee 30,000.\newline
Feb 10: Purchased goods for \rupee 40,000 on credit from M/s ABC Traders.\newline
Feb 12: Paid \rupee 5,000 for office rent in cash.\newline
Feb 17: Sold goods for \rupee 25,000 on credit to M/s XYZ Ltd.\newline
Feb 20: Paid \rupee 10,000 to M/s ABC Traders.\newline
Feb 22: Purchased office supplies for \rupee 2,000 in cash.\newline
Feb 24: Received \rupee 20,000 from M/s XYZ Ltd.\newline
Feb 27: Paid wages \rupee 4,000 in cash.\newline
Feb 28: Withdrew \rupee 3,000 for personal use.}{10}{3}{2}
\end{cietable}
\btfoot{|c|c|c|*{10}{C{0.85cm}|}}{\hline
\multirow{3}{*}{\bfseries Marks Distribution} & \multicolumn{2}{c|}{\bfseries Particulars} & CO1 & CO2 & CO3 & CO4 & CO5 & L1 & L2 & L3 & L4 & L5\\\cline{2-13}
 & \multirow{2}{*}{\shortstack{Max\\Marks}} & Quiz & 10 & -- & -- & -- & -- & 06 & 04 & -- & -- & --\\\cline{3-13}
 & & Test & 35 & 15 & -- & -- & -- & -- & 20 & 30 & -- & --\\\hline}
\stars{*****}
\end{iempaperB}

% ---------- IM255TBB Enterprise Information Systems (CIE - I)
\begin{iempaperB}
\paperinfo{ayear={2024-2025 (ODD Sem)},date={27\textsuperscript{th} November 2024},code={IM255TBB},sem={V},exam={CIE -- I},maxmarks={10 + 50},duration={120 Mins},
 title={ENTERPRISE INFORMATION SYSTEMS (Elective)}}
\iemheadB
\ptitle{PART -- A}
\begin{cietable}{M}{BT}{CO}
\q{1}{List two key features of a Customer Relationship Management (CRM) system.}{02}{L2}{CO1}
\q{2}{Mention two types of data commonly handled by EIS.}{02}{L2}{CO1}
\q{3}{Name two popular ERP software solutions.}{02}{L1}{CO2}
\q{4}{What is the primary function of Supply Chain Management (SCM) within EIS?}{02}{L1}{CO1}
\q{5}{What are the core modules of an ERP system?}{02}{L2}{CO2}
\end{cietable}
\ptitle{PART -- B}
\begin{cietable}{M}{BT}{CO}
\q{1.}{Analyse the benefits and challenges of implementing an EIS in a manufacturing organization. Support your answer with examples.}{10}{L2}{CO1}
\q{2.}{How can information technology support a company's business processes and decision making and give it competitive advantage? Give examples to illustrate your answer.}{10}{L3}{CO3}
\q{3.}{Highlight some of the issues you have faced in using SAP implemented in your institution. How you resolved those issues.}{10}{L3}{CO2}
\q{4.}{Evaluate the role of EIS in Customer Relationship Management (CRM) with specific examples of how companies have leveraged EIS for CRM.}{10}{L2}{CO3}
\q{5.}{XYZ Manufacturing Company has been facing inefficiencies in its operations due to a lack of integration between its departments. The procurement, inventory, production, and sales teams operate on standalone systems, leading to delays, data inaccuracies, and increased costs. To address these issues, the company has decided to implement an Enterprise Resource Planning (ERP) system. However, after six months of implementation, the company is still struggling with system adoption and is not seeing the expected improvements.\newline
Analyze the challenges XYZ Manufacturing is facing in implementing the ERP system. Suggest strategies to overcome these challenges and ensure successful adoption of the ERP system.}{10}{L4}{CO4}
\end{cietable}
\btfoot{|c|c|c|*{11}{C{0.82cm}|}}{\hline
\multirow{3}{*}{\bfseries Marks Distribution} & \multicolumn{2}{c|}{\bfseries Particulars} & CO1 & CO2 & CO3 & CO4 & CO5 & L1 & L2 & L3 & L4 & L5 & L6\\\cline{2-14}
 & \multirow{2}{*}{\shortstack{Max\\Marks}} & Quiz & 06 & 04 & -- & -- & -- & 04 & 06 & -- & -- & -- & --\\\cline{3-14}
 & & Test & 10 & 10 & 20 & 10 & -- & -- & 20 & 20 & 10 & -- & --\\\hline}
\stars{*****}
\end{iempaperB}

% ---------- IM355TBD Advanced Decision Modelling (CIE - 1)
\begin{iempaperB}
\paperinfo{ayear={2024-2025 (ODD Sem)},usn=0,date={27\textsuperscript{th} November 2024},code={IM355TBD},sem={VI},exam={CIE -- 1},maxmarks={50},duration={90 Min},
 title={ADVANCED DECISION MODELLING (Professional Core Elective)}}
\iemheadB
\begin{cietable}{M}{BT}{CO}
\qpart{Part A}
\q{1.}{What is the primary objective of sequencing problems?}{02}{1}{1}
\q{2.}{How does the concept of `setup time' impact sequencing problems?}{02}{3}{3}
\q{3.}{What is the significance of `priority rules' in sequencing problems?}{02}{2}{2}
\q{4.}{What is the difference between `flow shop' and `job shop' sequencing?}{02}{2}{2}
\q{5.}{How does the concept of `setup time' impact sequencing problems?}{02}{3}{3}
\qpart{Part B}
\q{1}{Explain the following terminologies:\newline
\hspace*{0.5cm}i. Idle time on machine\hspace{0.8cm}ii. Total elapsed time\hspace{0.8cm}iii. No passing rule\newline
\hspace*{0.5cm}iv. Number of machines\hspace{1.0cm}v. Processing order}{10}{1}{1}
\q{2}{Find the sequence that minimizes the total elapsed time (in hours) required to complete the following tasks on two machines.
\qtabc{|l|c|c|c|c|c|c|c|c|c|}{\hline \textit{Task} & \textit{A} & \textit{B} & \textit{C} & \textit{D} & \textit{E} & \textit{F} & \textit{G} & \textit{H} & \textit{I}\\\hline \textit{Machine I} & 2 & 5 & 4 & 9 & 6 & 8 & 7 & 5 & 4\\\hline \textit{Machine II} & 6 & 8 & 7 & 4 & 3 & 9 & 3 & 8 & 11\\\hline}}{10}{2}{2}
\q{3}{Determine the optimal sequence of jobs that minimize the total elapsed time, based on the following information. Processing time on machines is given in hours and passing is not allowed.
\qtabc{|l|c|c|c|c|c|c|c|}{\hline \textit{Job No} & \textit{A} & \textit{B} & \textit{C} & \textit{D} & \textit{E} & \textit{F} & \textit{G}\\\hline \textit{Machine $M_1$} & 3 & 8 & 7 & 4 & 9 & 8 & 7\\\hline \textit{Machine $M_2$} & 4 & 3 & 2 & 5 & 1 & 4 & 3\\\hline \textit{Machine $M_3$} & 6 & 7 & 5 & 11 & 5 & 6 & 12\\\hline}}{10}{3}{3}
\q{4}{There are four jobs each of which has to be processed on machines A, B, C, D, E and F in the order ABCDEF. Processing time in hours is given in the table below. Find the optimal sequencing of jobs, minimum time required to process these jobs and the idle time for each of the machines.
\qtabc{|c|c|c|c|c|c|c|}{\hline Job & A & B & C & D & E & F\\\hline 1 & 15 & 8 & 6 & 14 & 6 & 36\\\hline 2 & 17 & 7 & 9 & 10 & 15 & 22\\\hline 3 & 21 & 7 & 12 & 9 & 11 & 19\\\hline 4 & 18 & 6 & 11 & 12 & 14 & 17\\\hline}}{10}{2}{2}
\q{5}{Two major parts $P_1$ and $P_2$ for a product require processing through machine centres. The technological sequence of the parts on six machines and manufacturing times on each machine are
\qtabc{|l|l|c|c|c|c|c|c|}{\hline \multirow{2}{*}{Part $P_1$} & Machine sequence: & \textit{C} & \textit{A} & \textit{E} & \textit{F} & \textit{D} & \textit{B}\\\cline{2-8} & Time (hrs): & 2 & 3 & 4 & 5 & 6 & 1\\\hline \multirow{2}{*}{Part $P_2$} & Machine sequence: & \textit{B} & \textit{A} & \textit{E} & \textit{F} & \textit{C} & \textit{D}\\\cline{2-8} & Time (hrs): & 3 & 2 & 5 & 3 & 2 & 3\\\hline}
What would be the optimal scheduling to minimize the total processing time for these two parts? Find also the total elapsed time. For each machine specify the part that should be processed first.}{10}{3}{3}
\end{cietable}
\btfoot{|c|c|c|*{10}{C{0.85cm}|}}{\hline
\multirow{3}{*}{\bfseries Marks Distribution} & \multicolumn{2}{c|}{\bfseries Particulars} & CO1 & CO2 & CO3 & CO4 & L1 & L2 & L3 & L4 & L5 & L6\\\cline{2-13}
 & Quiz & Max Marks & 1 & 2 & 2 & -- & 1 & 2 & 2 & -- & -- & --\\\cline{2-13}
 & Test & Max Marks & 10 & 20 & 20 & -- & 10 & 20 & 20 & -- & -- & --\\\hline}
\stars{*******}
\end{iempaperB}

\end{document}
